\section*{5. 用“初等”方法计算对偶簇的次数}

(当 $\dim X$ 为偶数\underline{或} $\operatorname{char}(k)\ne2$ 时的另一种方法参见 XVIII, 3.2.)

继续使用 2.2 的记号: 因此 $X\hookrightarrow\mathbb P^r$ 为维数 $n$ 的光滑不可约子簇.

\subsubsection*{5.1.}
首先考虑态射 (3.1.1)
\begin{equation}\tag{5.1.0}
\nu:\mathbb P(N)\longrightarrow X\times\check{\mathbb P}^{r}.
\end{equation}
由定义, 余法层 $\check N$ 出现在 $X$ 上凝聚层的正合列
\begin{equation}\tag{5.1.1}
0\longrightarrow\check N
\longrightarrow\Omega^1_{\mathbb P^r/k}|_X
\longrightarrow\Omega^1_X
\longrightarrow0.
\end{equation}
中. 另一方面, 用齐次坐标表示 $\mathbb P^r$ 上的微分形式, 得到 $\mathbb P^r$ 上层的正合列
\begin{equation}\tag{5.1.2}
0\longrightarrow\Omega^1_{\mathbb P^r/k}
\longrightarrow\mathcal O_{\mathbb P^r}(-1)^{r+1}
\longrightarrow\mathcal O_{\mathbb P^r}
\longrightarrow0.
\end{equation}
因而其在 $X$ 上的限制为正合列
\begin{equation}\tag{5.1.3}
0\longrightarrow\Omega^1_{\mathbb P^r/k}|_X
\longrightarrow\mathcal O_X(-1)^{r+1}
\longrightarrow\mathcal O_X
\longrightarrow0.
\end{equation}
\footnote{印本在将此列限制到 $X$ 后, 最后一项仍保留为 $\mathcal O_{\mathbb P^r}$; 限制后的项为 $\mathcal O_X$.}
结合 (5.1.1) 和 (5.1.3), 得到包含
\begin{equation}\tag{5.1.4}
\check N\lhook\joinrel\longrightarrow\mathcal O_X(-1)^{r+1},
\end{equation}
转置后得到满射
\begin{equation}\tag{5.1.5}
\mathcal O_X(1)^{r+1}\longrightarrow N.
\end{equation}
因而有闭浸入
\begin{equation}\tag{5.1.6}
\alpha:\mathbb P(N)
\longrightarrow\mathbb P\bigl(\mathcal O_X(1)^{r+1}\bigr)
\simeq X\times\check{\mathbb P}^{r},
\end{equation}
这恰好就是 (5.1.0).

